%% Die Sprachdatei von helsinki für en-GB (#377). Was eine Sprachdatei ist und
%% welche Regeln für sie gelten, steht im Kopf von helsinki-de-de.tex.

%% ---------------------------------------------------------------- Zahlen
%% Dezimalpunkt und Komma als Tausendertrenner (ADR-0021). Ab wie vielen
%% Stellen gruppiert wird, wählt die Vorlage.
\sisetup{
  output-decimal-marker = {.},
  group-separator       = {,},
}
%% Das €-Zeichen steht britisch vor der Zahl und ohne Abstand, das Minus
%% davor: −€200.00. siunitx setzt kein Zeichen vor die Zahl, deshalb trennt
%% \betrag@zeichen das Vorzeichen ab und setzt es selbst, als \textminus wie
%% siunitx (#320). Das erste Zeichen genügt: Kotlin liefert die Ziffern mit
%% Punkt und einem Minus davor, nie ein Plus oder Leerraum (ADR-0021).
%%
%% Das @ im Namen macht die Hilfe zu einem Makro allein dieser Datei; die
%% übrigen Sprachdateien brauchen sie nicht zu kennen.
\makeatletter
\newcommand{\betrag}[1]{\betrag@zeichen#1\relax}
\def\betrag@zeichen#1#2\relax{\ifx-#1\textminus€\num{#2}\else€\num{#1#2}\fi}
\makeatother
%% Das Prozentzeichen ohne Abstand: 19%.
\newcommand{\prozent}[1]{\num{#1}\%}

%% ------------------------------------------------------------- Seiten
\newcommand{\wortseitevon}[2]{Page #1 of #2}
\newcommand{\wortfortsetzung}{Continued on the next page}

%% ------------------------------------------------------ Kopf und Anschrift
\newcommand{\wortgeschaeftsfuehrung}{Managing director}
\newcommand{\wortustid}{VAT ID}
\newcommand{\wortkundennr}{Customer no.:}
\newcommand{\wortauftragsnr}{Our order no.:}
\newcommand{\wortvereinnahmtam}{Payment received on}
\newcommand{\wortreferenz}{Your reference:}
%% "Due in full" für "ohne Abzug": Es sagt dasselbe, ohne dass jemand nach
%% einem Abzug fragt.
\newcommand{\wortzahlbarbis}{Due in full by}

%% ----------------------------------------------------- Positionen, Summen
\newcommand{\wortbeschreibung}{Description}
\newcommand{\wortmenge}{Quantity}
\newcommand{\worteinzelpreis}{Unit price}
\newcommand{\wortrabatt}{Discount}
\newcommand{\wortbetrag}{Amount}
\newcommand{\wortnettobetrag}{Net amount}
%% Satz und, bei mehreren Steuergruppen, ihre Grundlage – sonst ist sie leer.
\newcommand{\wortsteuer}[2]{plus #1 VAT\ifblank{#2}{}{ on #2}}
\newcommand{\wortgesamtbetrag}{Total amount}

%% --------------------------------------------------------- Schluss und Fuß
\newcommand{\wortgruss}{Kind regards}
\newcommand{\wortbankverbindung}{Bank details}
\newcommand{\wortkontoinhaber}{Account holder}
\newcommand{\wortinstitut}{Bank}
\newcommand{\wortregistergericht}{Registry court}
